प्रथम-क्रम भाषेची \emph{चिन्हार्थमीमांसा} त्या भाषेसाठीच्या \emph{रचनेद्वारे} दिली जाते. रचनेत एक \emph{प्रांत} असतो आणि प्रत्येक स्थिरचिन्हाला त्या प्रांतातील एक घटक नेमला जातो. फलनचिन्हांचा अर्थ प्रांतावरील फलनांनी, तर संबंधचिन्हांचा अर्थ प्रांतावरील संबंधांनी ठरवला जातो. चरांच्या संचापासून प्रांतात जाणारे फलन म्हणजे \emph{चर-मूल्यांकन}. \emph{पूर्ति} संबंध रचना, चर-मूल्यांकने आणि सूत्रे यांना जोडतो. $\Sat{M}{!A}[s]$ हे $!A$ च्या रचनेवरील विगमनाने परिभाषित केले जाते. $\Sat{M}{!A}[s]$ हे रचनेच्या प्रांतावर, $!A$ मध्ये प्रत्यक्ष येणाऱ्या अतार्किक चिन्हांच्या अर्थनिर्धारणावर आणि $!A$ मधील मुक्त चरांना~$s$ ने दिलेल्या मूल्यांवर अवलंबून असते. विशेषतः, $!A$ मध्ये मुक्त चर नसतील तर ते $s$ वर अवलंबून नसते. म्हणून $!A$ हे वाक्य असेल, तर कोणत्याही एका (किंवा सर्व)~$s$ साठी $\Sat{M}{!A}[s]$ असेल तेव्हा $\Sat{M}{!A}$ असे म्हणतो.

पूर्ति-संबंध हा सर्व चिन्हार्थविषयक संकल्पनांचा आधार आहे. प्रत्येक रचनेत ज्या वाक्याची पूर्ति होते ते वाक्य \emph{वैध} असते, म्हणजे $\Sat{{}}{!A}$. वाक्यसंच~$\Gamma$ मधील प्रत्येक वाक्याची पूर्ति करणाऱ्या सर्व $\Struct{M}$ साठी $\Sat{M}{!A}$ असेल तेव्हा आणि केवळ तेव्हाच वाक्य $!A$ त्या संच~$\Gamma$ पासून \emph{तार्किकरीत्या निष्पन्न} होते, म्हणजे $\Gamma \Entails !A$. संच~$\Gamma$ मधील प्रत्येक वाक्याची पूर्ति करणारी एखादी रचना असेल तेव्हा आणि केवळ तेव्हाच तोच संच~$\Gamma$ \emph{पूर्ततायोग्य} असतो, अन्यथा तो \emph{अपूर्ततायोग्य} असतो. या संकल्पना परस्परांशी संबंधित आहेत. उदाहरणार्थ, $\Gamma \Entails !A$ तेव्हा आणि केवळ तेव्हाच जेव्हा $\Gamma \cup \{\lnot !A\}$ अपूर्ततायोग्य असतो.
